% !TeX program = lualatex
\documentclass[paper=a4, fleqn]{jlreq}

% ===== 日本語・数式 =====
\usepackage{luatexja}
\usepackage{amsmath,amssymb}
\usepackage{ascmac}

% ===== レイアウト =====
\usepackage[
	top=18mm,
	bottom=20mm,
	left=9mm,
	right=9mm
]{geometry}

% ===== 図・装飾 =====
\usepackage{tikz}
\usepackage{multicol}
\usepackage{tcolorbox}
\usepackage{graphicx}
\usepackage{xcolor}

% ===== 既存教材用 =====


	\tcbuselibrary{skins}
\usetikzlibrary{patterns}

% ===== 各種設定 =====
\renewcommand{\labelenumi}{(\theenumi)\ }
\newcommand{\ctext}[1]{%
	\raise0.2ex\hbox{\textcircled{\scriptsize #1}}%
}

\newtcolorbox{mysimplebox}[1]{%
	colframe=black,
	colback=white,
	coltitle=black,
	colbacktitle=white,
	boxrule=0.8pt,
	arc=0mm,
	fonttitle=\sffamily\bfseries,
	enhanced,
	attach boxed title to top left={
		xshift=3mm,
		yshift=-4mm
	},
	boxed title style={frame hidden},
	title=#1
}

\setlength{\columnseprule}{0.8pt}

\begin{document}

% ==================================================
% 表紙
% ==================================================

\begin{center}
\mbox{}\\[18mm]

\huge
{\textbf{数\qquad\quad 学}}\\[24mm]

\small
	{\textbf{〈監督者の指示があるまで開いてはいけない〉}}\\
\end{center}

\vspace{9mm}

\small

\qquad\qquad\qquad\qquad\qquad\qquad
1.\quad
試験開始後，まず解答用紙に自分の受験番号と氏名を正しく記入しなさい。\\[-2mm]

\qquad\qquad\qquad\qquad\qquad\qquad
2.\quad
試験開始後，速やかに問題冊子に落丁や乱丁がないか確認しなさい。\\[-2mm]

\qquad\qquad\qquad\qquad\qquad\qquad
\phantom{2.\;\;}
落丁や乱丁があった場合は，手を挙げなさい。\\[-2mm]

\qquad\qquad\qquad\qquad\qquad\qquad
3.\quad
解答用紙に印刷されていない問いの番号は各自で記入しなさい。\\[-2mm]

\qquad\qquad\qquad\qquad\qquad\qquad
4.\quad
下書きは問題用紙の余白を利用しなさい。\\[-2mm]

\qquad\qquad\qquad\qquad\qquad\qquad
5.\quad
問題用紙は試験終了後，持ち帰ってもよい。\\[-2mm]

\qquad\qquad\qquad\qquad\qquad\qquad
\phantom{5.\;\;}
ただし，試験途中には持ち出してはいけない。

% ==================================================
% 問題ページ
% ==================================================

\newpage
\mbox{}\\[6mm]

\qquad\qquad\qquad
{\LARGE\textbf{１}}.\;
\normalsize
次の \fbox{\rule{0pt}{2.0ex}\rule{0.7cm}{0pt}}
にあてはまる適切な数値を解答欄に記入せよ。\\[1.5mm]

\qquad\qquad\quad\;\;
2つの袋A，Bがあり，袋Aには赤玉2個，白玉1個が，袋Bには赤玉1個，白玉2個が入\\
\qquad\qquad\quad\;\;
っている。袋A，Bから玉を1個ずつ取り出し，取り出した2個の玉の色が異なるときは，そ\\
\qquad\qquad\quad\;\;
の2個の玉を互いに入れ替えて元の袋に戻し，同じ色であるときは，そのまま元の袋に戻すと\\
\qquad\qquad\quad\;\;
いう操作を繰り返し行う。このとき，\\[-1mm]

\qquad\qquad\quad\;\;
・2回目の操作を終えたとき，袋Aの中に残っている赤玉が1個である確率は
\;\fbox{(ア)}\;,\\[-2mm]

\qquad\qquad\quad\;\;
・4回目の操作を終えたとき，袋Aの中に残っている赤玉が0個である確率は
\;\fbox{(イ)}\;\\

\qquad\qquad\quad\;\;
である。
\end{document}
